\chapter{Discussion}
\label{ch:discussion}

\section{Interpreting the coverage results}

The main empirical finding of this dissertation is that standard split conformal prediction over-covers easy regions and badly under-covers difficult regions, while the proposed adaptive method keeps coverage much closer to nominal 90\% across all difficulty levels (\cref{tab:tertile,tab:sixbin}). The same qualitative pattern was reported by \citet{amirishahbazi2026} for both of their case studies, despite the difference in problem structure discussed in \cref{sec:adaptation}. Therefore, the similarity suggests that the adaptive approach may work across different types of prediction problems. The key mechanism driving the improvement is the normalisation of nonconformity by a locally informative difficulty estimate, which can improve the allocation of interval width across different levels of difficulty. This is an empirical improvement rather than a formal theoretical guarantee of conditional coverage, since the formal conformal guarantee is for marginal coverage under the required exchangeability assumptions. The exact geometric construction is different, with previous work using geodesic caps for a manifold-valued response, while this dissertation is on scalar prediction intervals whose widths are determined using a geodesic-neighbourhood difficulty estimator. Despite this difference, the important component is whether the difficulty estimator is sufficiently informative. Therefore, the particular geometric implementation may be less important than obtaining a useful local estimate of prediction difficulty.
 
That said, the two settings are not equally challenging for the mechanism to succeed in. In this current dissertation setting, the covariate is the exact geographical location of each city. Coordinates are known directly and the geodesic $k$-NN neighbourhood used for $\hat\sigma(x)$ can therefore be computed directly on $S^2$. There is no additional uncertainty in determining the neighbourhood apart from numerical and geometric calculations. In \citeauthor{amirishahbazi2026}'s setting, both the covariate and response can involve noise or model-dependent uncertainty. The response is also manifold-valued, so $\hat\sigma(x)$ can be more complicated as errors from the base predictor may interact with the geometry of the manifold-valued response. 

The higher $r$ does not necessarily mean the proposed method is universally better. This gap between the correlation obtained here ($r = 0.634$) and \citeauthor{amirishahbazi2026}'s real-data case study ($r = 0.516$) has at least two plausible, non-exclusive explanations. The first is the setting difference just described: exact spatial coordinates, a scalar response, and direct geodesic neighbourhood calculation avoid the interaction between predictor error and output curvature present in \citeauthor{amirishahbazi2026}'s manifold-response setting. A second, non-exclusive explanation is that near-surface temperature is strongly spatially autocorrelated at the scale of tens to hundreds of kilometres, which may give a $k$-nearest-neighbour residual estimator a genuinely easier signal to exploit than for geomagnetic field anomalies, which may vary more locally. The present study cannot fully separate these two explanations, as doing so would require applying both methods to a common underlying field. The relative stability of $r$ across spherical harmonic degree ($r \approx 0.61$-$0.68$ for $L = 8$-$30$, \cref{tab:degree-robustness}) is weak but suggestive evidence that the correlation reflects the spatial structure of the temperature field rather than any single model configuration, though this is not conclusive. Therefore, the two correlation values should be compared as empirical results from different problem settings rather than as a direct ranking of the two methods.
 
\section{Why interval width is comparable, not just coverage}

The adaptive method's improvement is captured clearly in \cref{tab:main-coverage} where it improves coverage without producing wider intervals overall and its mean width is slightly narrower on average than the standard method ($4.59^\circ\mathrm{C}$ versus $4.65^\circ\mathrm{C}$). This is possible because the adaptive method redistributes width instead of just uniformly increasing it. It narrows intervals for easy-tertile locations while widening them for hard-tertile locations. Standard conformal uses a common interval width, even though prediction difficulty varies spatially. This means easy locations may receive intervals that are unnecessarily wide and difficult locations may receive intervals that are not wide enough. Adaptive conformal provides a more appropriate allocation of uncertainty. As a result, the overall mean interval width remains approximately the same, while producing a large improvement in the conditional-coverage metrics that the mean alone cannot capture. This is the same point that \citeauthor{amirishahbazi2026} make about their naive-coordinate baseline wasting prediction-region area rather than simply causing incorrect coverage. In the current scalar-response setting, a similar idea applies where uniform interval width does not account for different prediction difficulties. Therefore, the advantage of the adaptive method is not simply higher coverage. It provides more balanced conditional coverage, similar or slightly smaller average interval width, narrower intervals for easier locations, and wider intervals for genuinely difficult locations. This suggests that the adaptive method provides a more efficient allocation of uncertainty across the spatial domain.
 
\section{Computational considerations}

The pipeline's main computational costs in this dissertation come from three parts: (i) the five-fold cross-validated hyperparameter grid search over $21$ values of $\lambda$ and $5$ values of the Sobolev order $k$ at $L=26$, which requires $5\times21\times5=525$ fold-level model evaluations; (ii) the five-fold cross-validation used to construct the out-of-fold residual field for $\hat\sigma(x)$, which requires five additional model refits; and (iii) the degree-robustness experiment, which repeats these procedures at eight different values of $L$. For all repeated grid-search calculations, PyTorch is used for spherical-harmonic matrix operations, matrix multiplication, and regularised linear-system solving, while SciPy is used for associated Legendre-function evaluations performed on the CPU. Within each fold, the design matrix, Gram matrix $A^HA$, and right-hand side $A^H\hat f$ are computed once and reused across the different $(\lambda,k)$ combinations. This avoids repeatedly constructing the same matrices. The final model fit is performed separately using the NumPy/SciPy implementation. As the tested degree becomes larger, such as $L=30$, the size of the normal-equation system also increases, with dimension $(L+1)^2=961$. Compared with $729$ at $L=26$, this results in a larger linear system and higher computational cost. The experiments were feasible on the available computer hardware, although increasing the degree, substantially refining the hyperparameter grid, testing more Sobolev orders, or using more CV folds would increase the computational cost. However, once training and calibration are complete, evaluating the calibrated model at a new query point is comparatively inexpensive. This is because it only requires evaluation of the spherical-harmonic basis, estimation of the local difficulty from the precomputed residual field, and computation of the corresponding prediction interval. Overall, the training and model-selection process is more computationally expensive than prediction at a new point.
 
\section{Comparison with alternative spatial uncertainty quantification approaches}

It is worth investigating and comparing the conformal approach against the geostatistical alternative mentioned in \cref{ch:intro}, Gaussian Process/Kriging, a natural alternative from spatial statistics. A Gaussian Process produces a probabilistic prior over the spatial field, which requires a specifying covariance/kernel function (e.g. spherical Matérn). This kernel is used to represent assumptions about spatial correlation and smoothness. A Gaussian Process will therefore produce a full posterior predictive distribution, such as a predictive mean, variance and intervals directly. This is much richer than the conformal interval, but its uncertainty estimates depend heavily on the chosen kernel and noise model. So if the assumed covariance structure is wrong or misleading, the predictive variance may be poorly calibrated. This might potentially affect my data because of spatially varying/terrain-related temperature patterns documented in \cref{sec:geography}. The resulting posterior variance can be easily systematically miscalibrated with no finite-sample mechanism to correct it. In contrast, the conformal guarantee of \cref{sec:cp-background} holds regardless of model misspecification. In conclusion, a Gaussian Process provides richer probabilistic uncertainty, but its uncertainty is dependent on correctly specifying the spatial covariance structure. 

My approach takes a different route. It keeps the existing spherical harmonic reconstruction model and applies conformal prediction on top of that predictor. It does not require specifying a Gaussian Process covariance kernel. In my case, calibration residuals are used to determine the prediction interval. Therefore, the uncertainty calibration does not depend on the spherical harmonic model having a correctly specified probabilistic covariance structure. The conformal coverage guarantee is model-agnostic, but not assumption-free. The advantage is therefore reliable interval calibration without replacing or rebuilding the underlying spatial predictor. Conformal prediction sacrifices the full probabilistic output of a Gaussian Process, but provides a model-agnostic way to calibrate prediction intervals on the existing predictor.

Gaussian process regression is also computationally more demanding because its covariance matrix has dimension $n\times n$, with a cubic computational cost of $O(n^3)$ for a direct factorisation. In the present dataset, with $n=2{,}106$ training points, it is feasible to compute but the cost grows quickly with more observations. On the other hand, the spherical harmonic formulation works in a coefficient space with dimension $(L+1)^2$. Dense factorisation therefore has a worst-case complexity of $O(((L+1)^2)^3)=O((L+1)^6)$. This distinction is important because the two approaches scale with different quantities. The Gaussian Process approach is driven primarily by the number of observations, whereas the spherical harmonic calculation is driven primarily by the number of basis functions. For the present setting, where $L\ll n$, the latter provides a substantially smaller linear system. It should be noted that theoretical complexity does not necessarily equal actual runtime, because implementation, hardware, matrix structure, and approximation methods also matter. The conformal approach adopted in this dissertation therefore trades the fully probabilistic output of a Gaussian Process for a model-agnostic marginal coverage guarantee that can be applied to the existing spherical harmonic predictor without redesigning the underlying reconstruction model.

\section{Broader implications}

Going back to the practical motivation from the introduction \cref{ch:intro}, a downstream user is also expected to know where the reconstructed field can be trusted, as a single global accuracy measure is not enough. The tertile results in \cref{tab:tertile} show that standard conformal has a nominal coverage of 90\%, but the actual coverage in the highest-difficulty tertile is only 76.7\%. This is a substantial difference from the target coverage, with the under-coverage mainly concentrated in particular spatial regions across the dataset. The geographic analysis in \cref{sec:geography} suggests some possible physical and geographical reasons behind these patterns. A user may assume that a nominal 90\% interval provides roughly 90\% reliability everywhere, but the results show that coverage can vary across different levels of prediction difficulty. The greatest risk occurs in areas where the predictor is most uncertain. Therefore, the adaptive method improves the numerical results while also providing more informative uncertainty estimates for downstream users. For example, a regional climate analyst may need to decide whether the reconstructed field over the Bolivian Altiplano is reliable enough for further modelling. A decision-maker could also use this information to identify areas where additional monitoring could improve confidence in the reconstructed field. The spatial uncertainty map \cref{fig:opened-globe} helps to show how the model performs across different regions, providing information that is not captured by an overall accuracy metric.

Adaptive conformal prediction does not address every type of spatial uncertainty. The conformal guarantee applies to the data distribution represented by the calibration and test data. It relies on the relevant exchangeability assumption and does not extend to arbitrary unseen regions. For instance, open-ocean locations are outside the spatial regime represented by the training data, as the available observations are limited to land-based locations (\cref{ch:data}). Therefore, oceanic locations differ from the regions represented in the observed data. This makes the issue one of extrapolation and generalisation, rather than just calibration.

A narrow interval at an oceanic location may appear reassuring, but it does not necessarily mean that the prediction is accurate there. The conformal calibration does not validate performance in a regime that is absent from the calibration data. In \cref{fig:opened-globe}, the uncertainty appears visually unremarkable over some oceanic grid points. However, the current analysis does not establish why this occurs, so it should not be interpreted as evidence that the method underestimates oceanic uncertainty. It is important to distinguish between calibration and accuracy. Calibration refers to how reliable the prediction intervals are under the represented data distribution, while accuracy refers to how well the underlying prediction performs in a new or structurally different region. Therefore, a calibrated interval should not be interpreted as evidence that the model is accurate everywhere, especially in regions that are poorly represented or absent from the observations. As a result, users should interpret uncertainty maps together with the spatial coverage of the training and calibration data. Regions outside the observed data regime require additional validation or physical observations. The proposed method improves uncertainty calibration within the represented data distribution, but it does not eliminate the problem of extrapolation into regions that are not represented in the observations. This distinction is important when using the reconstruction and uncertainty maps in practice, as good conformal calibration does not mean that the model is accurate everywhere.

\section{Limitations}

Several additional limitations should be acknowledged. 

First, the underlying dataset itself is a spatial and temporal average (Berkeley Earth's monthly city-level average), rather than a direct physical measurement, so it is not appropriate to interpret the residuals as measurement error. The reconstruction and its uncertainty should be described as a model fit to this processed dataset, rather than to the underlying physical measurement itself. This issue was already mentioned for the point-estimate model earlier and is carried over without change by this uncertainty-quantification layer. 

Second, the difficulty estimator's cross-validated residual field is estimated using only $2{,}106$ training points, and its own accuracy does not remain perfectly stable as model complexity changes, as shown in \cref{fig:degree-robustness}. A larger or more geographically diverse dataset could potentially improve the estimation of $\hat\sigma(x)$ further. 

Third, the specific choice $k_\sigma=20$ for the difficulty-estimator neighbourhood size was taken directly from \citet{amirishahbazi2026}. No independent tuning process was conducted for this dataset; therefore, no sensitivity analysis over $k_\sigma$ was conducted.

Fourth, the adaptation described in \cref{sec:adaptation} illustrates the adaptation motivated by the structural mismatch between a scalar-response problem and a manifold-response framework. It should not be seen as a validated theoretical result similar to the theoretical result in \citeauthor{amirishahbazi2026}, which is already well established and proved for the manifold-response case. The coverage guarantee obtained here instead follows from the standard (Euclidean, locally-adaptive) split-conformal argument of \citet{lei2018} and \citet{romano2019}. The geodesic input-neighbourhood construction might be one possible choice of difficulty estimator when it is geometrically well motivated.

Fifth, the reported test-set metrics in \cref{sec:main-results,tab:sixbin} are computed from a single fixed calibration and test split. This is the same split used throughout \cref{ch:results}, rather than being averaged over repeated random splits as in \citeauthor{amirishahbazi2026}'s $300$-trial and $100$-trial Monte Carlo evaluations. For the current dataset, a single split is already enough to demonstrate the conditional miscalibration of the standard method and the improvement obtained with the adaptive method. This also matches how the base predictor itself was evaluated in the earlier dissertation. The coverage and interval-width results are based on one fixed calibration and test split. Both calibration and test sets contain $n_{cal}=n_{te}=351$. Therefore, the exact numerical values have some unavoidable sampling variability, and the reported 88\%, 90\%, width, etc. should not be interpreted as exact population-level quantities. The main qualitative pattern is repeated across all eight values of $L$, as shown in \cref{tab:degree-robustness}. Each degree involves separate model fitting, calibration, and evaluation. This provides some reassurance that the main finding is not specific to one particular model degree. However, the eight experiments do not use independent random data splits. They are based on the same underlying partition. Therefore, they should not be described as eight independent replications. They mainly demonstrate that the result is reasonably consistent when model complexity changes. This still supports the robustness of the qualitative pattern to the choice of $L$, but does not quantify how much the results would change under a different random train/calibration/test partition. Therefore, it provides weaker evidence than repeated random-split experiments. A stronger evaluation would repeat the complete experiment over multiple random splits. This would allow examination of variation in empirical coverage and mean interval width, and stability of the adaptive-vs-standard comparison. It could also report mean, standard deviation, and confidence intervals. This would provide uncertainty estimates for the performance metrics themselves. Compared with the prior work of \citeauthor{amirishahbazi2026}, which used repeated Monte Carlo trials, a similar repeated-trial evaluation would improve this experimental design. It would also provide stronger evidence for the generality of the numerical results. 
 
% ============================================================
